\documentclass[sans-serif]{sfs-2487-2024}

% amssymb tarjoaa tehtävälistan valintaruudut ($\square$, $\boxtimes$),
% samat kuin pandoc käyttää markdown-vastineen ''- [ ]'' / ''- [x]'' -kohdissa.
\usepackage{amssymb}

% Käyttöohje, jossa on kuvatekstillisiä kuvia tekstin tasossa ja tekstin
% sisennyskohdassa (6.5.2), numeroituja vaiheluetteloita sisäkkäisine
% luetteloineen (6.5.3), tehtävälista valintaruuduin sekä korostuksia
% (lihavointi, kursiivi, tasavälinen koodimerkintä). Nämä vastaavat
% markdown-version ../markdown/esimerkki-kayttoohje.md merkintöjä. Vaihtoehto
% [sans-serif] latoo asiakirjan groteskilla, joka vastaa standardin omien
% esimerkkien ulkoasua.

\logo{\includegraphics{logo-organisaatio}}
\doctype{Käyttöohje}
\date{20.5.2024}
\author{Virve Virtanen}

\subject{Graafinen ohjeisto}
\title{Logon käyttö asiakirjoissa}

\contactinfo{Organisaatio Oy}{%
  Katuosoite\\
  Postinumero Postitoimipaikka\\
  viestinta@organisaatio.fi}

\begin{document}

\maketitle

\section{Logon perusversio}

Logon perusversiota käytetään \textbf{kaikissa} organisaation
asiakirjoissa. Logo sijoitetaan asiakirjan vasempaan yläkulmaan, ja
\emph{sähköisessä} asiakirjassa sille annetaan vaihtoehtoinen teksti,
esimerkiksi ''Organisaatio Oy:n logo''. Perusversion tiedostonimi on
\texttt{logo-organisaatio.pdf}.

\noindent\includegraphics{logo-organisaatio}
\captionof{figure}{Logon perusversio}

\section{Suoja-alue}

Logon ympärille jätetään aina vapaata tilaa vähintään logon
merkkiosan korkeuden~(x) verran. Suoja-alueelle ei sijoiteta
tekstiä eikä muita graafisia elementtejä.

\noindent\includegraphics{logo-suoja-alue}
\captionof{figure}{Logon suoja-alue, jonka leveys on merkkiosan korkeus~x}

\section{Logon lisääminen asiakirjaan}

Logo lisätään asiakirjamalliin seuraavasti:

\begin{enumerate}
  \item Avaa asiakirjamalli ja siirry asiakirjan alkuun.
  \item Lisää logotiedosto asiakirjan vasempaan yläkulmaan. Käytä
    tarkoitukseen sopivaa tiedostomuotoa:
    \begin{itemize}
      \item painotuotteissa \texttt{.pdf} tai \texttt{.eps}
      \item verkkojulkaisuissa \texttt{.png}
    \end{itemize}
  \item Kirjoita logolle vaihtoehtoinen teksti.
  \item Tarkista, että perusmetatietojen alue ei mene logon päälle.
\end{enumerate}

Jos logon koko tai sijoittelu sitä vaatii, perus- ja lisämetatietojen
alueita voi siirtää alemmas tai enemmän oikealle.

\section{Tarkistuslista ennen julkaisua}

\begin{itemize}
  \item[$\boxtimes$] Logo on asiakirjan vasemmassa yläkulmassa.
  \item[$\boxtimes$] Logolla on vaihtoehtoinen teksti.
  \item[$\square$] Suoja-alue on tarkistettu.
\end{itemize}

\forinformation{Viestintäosasto}

\makecontactinfo

\end{document}
